\documentclass[a4paper,12pt]{article}
\usepackage[utf8]{inputenc}
\usepackage[ngerman]{babel}
\usepackage{amsmath,amssymb}
\usepackage{geometry}\geometry{margin=2.5cm}
\usepackage{enumitem}
\usepackage{booktabs}
\usepackage{tcolorbox}
\tcbuselibrary{breakable}
\usepackage{xcolor}

\newtcolorbox{merkbox}[1][]{colback=blue!5, colframe=blue!40, fonttitle=\bfseries, title={Merke}, breakable, #1}
\newtcolorbox{analogie}[1][]{colback=green!5, colframe=green!40, fonttitle=\bfseries, title={Analogie}, breakable, #1}
\newtcolorbox{begriffe}[1][]{colback=yellow!10, colframe=yellow!60!black, fonttitle=\bfseries, title={Was bedeuten die Begriffe?}, breakable, #1}
\newtcolorbox{wastun}[1][]{colback=orange!5, colframe=orange!60!black, fonttitle=\bfseries, title={Was ist zu tun?}, breakable, #1}

\title{Erkl\"arung -- HW09: Lineare Abbildungen, Kern \& Bild}
\author{}
\date{}

\begin{document}
\maketitle

%=============================================================================
\part*{Teil 1: Grundlagen -- Was du wissen musst}

%-----------------------------------------------------------------------------
\section*{1. Was ist eine lineare Abbildung?}

Eine \textbf{lineare Abbildung} $T \colon V \to W$ ist eine Funktion zwischen Vektorr\"aumen, die zwei Eigenschaften erf\"ullt:
\begin{enumerate}[nosep]
    \item $T(u + v) = T(u) + T(v)$ (Addition wird respektiert).
    \item $T(\lambda v) = \lambda \, T(v)$ (Skalarmultiplikation wird respektiert).
\end{enumerate}

\begin{analogie}
Eine lineare Abbildung ist wie ein \textbf{strukturerhaltender \"Ubersetzer}: Wenn ich zwei Vektoren erst addiere und dann \"ubersetze, kommt dasselbe heraus, wie wenn ich sie erst einzeln \"ubersetze und dann addiere. Beispiel: ``mal 2 nehmen'' ist linear, ``$+5$'' ist es nicht (denn $u + v + 5 \neq (u+5) + (v+5)$).
\end{analogie}

\begin{merkbox}[title={Darstellungsmatrix}]
Jede lineare Abbildung $T \colon \mathbb{R}^n \to \mathbb{R}^m$ l\"asst sich durch eine $m \times n$-Matrix $A$ darstellen: $T(v) = A v$.

\medskip
\textbf{Spalten von $A$} = Bilder der Einheitsvektoren: $A = (T(e_1) \mid T(e_2) \mid \dots \mid T(e_n))$.
\end{merkbox}

%-----------------------------------------------------------------------------
\section*{2. Kern und Bild}

\begin{merkbox}[title={Definitionen}]
\begin{itemize}[nosep]
    \item \textbf{Kern}: $\text{Ker}(T) = \{v \in V : T(v) = 0\}$. ``Was wird auf null abgebildet?''
    \item \textbf{Bild}: $\text{Im}(T) = \{T(v) : v \in V\}$. ``Was kann \"uberhaupt herauskommen?''
\end{itemize}
Beides sind Unterr\"aume -- $\text{Ker}(T) \subseteq V$ (Definitionsbereich), $\text{Im}(T) \subseteq W$ (Zielraum).
\end{merkbox}

\begin{analogie}
$T$ ist eine ``Maschine''. Der \textbf{Kern} ist die Menge der ``vergeudeten Eing\"ange'' (kommen als $0$ raus, also Information verloren). Das \textbf{Bild} ist die Menge aller m\"oglichen Ausg\"ange.
\end{analogie}

%-----------------------------------------------------------------------------
\section*{3. Der Dimensionssatz (Rang-Nullit\"ats-Satz)}

\begin{merkbox}[title={Der wichtigste Satz}]
F\"ur jede lineare Abbildung $T \colon V \to W$ mit endlicher Definitionsdimension:
\[
\dim \text{Ker}(T) + \dim \text{Im}(T) = \dim V.
\]
\end{merkbox}

\begin{analogie}
Wenn die Maschine $n$ Eingangsregler hat ($\dim V = n$), dann gilt: ``vergeudete Regler'' (Kern) + ``wirksame Regler'' (Bild) $= n$. Jeder Regler ist entweder das eine oder das andere -- nichts wird doppelt gez\"ahlt.
\end{analogie}

%-----------------------------------------------------------------------------
\section*{4. Injektiv, Surjektiv, Bijektiv -- Verbindung zu Kern \& Bild}

\begin{merkbox}[title={Charakterisierungen}]
F\"ur $T \colon V \to W$:
\[
\textbf{injektiv} \iff \text{Ker}(T) = \{0\} \iff \dim \text{Ker}(T) = 0.
\]
\[
\textbf{surjektiv} \iff \text{Im}(T) = W \iff \dim \text{Im}(T) = \dim W.
\]
\[
\textbf{bijektiv} \iff \text{injektiv und surjektiv} \iff T \text{ invertierbar}.
\]
\end{merkbox}

\begin{merkbox}[title={Wichtige Konsequenz}]
F\"ur $T \colon V \to W$ mit $\dim V = \dim W$ (also Quadratmatrix):
\[
\text{injektiv} \iff \text{surjektiv} \iff \text{bijektiv}.
\]
\textbf{Aber:} wenn $\dim V \neq \dim W$, dann kann $T$ \textbf{niemals bijektiv} sein.
\begin{itemize}[nosep]
    \item $\dim V > \dim W$ (z.\,B.\ $\mathbb{R}^5 \to \mathbb{R}^3$): $T$ kann nicht injektiv sein (Kern hat mindestens Dimension $\dim V - \dim W$).
    \item $\dim V < \dim W$ (z.\,B.\ $\mathbb{R}^3 \to \mathbb{R}^5$): $T$ kann nicht surjektiv sein.
\end{itemize}
\end{merkbox}

%-----------------------------------------------------------------------------
\section*{5. Strategie -- wie l\"ose ich solche Aufgaben?}

\begin{merkbox}[title={Kochrezept}]
\begin{enumerate}[nosep]
    \item \textbf{Darstellungsmatrix $A$ aufstellen} (Spalten = $T(e_i)$).
    \item \textbf{Gau\ss-Elimination} auf $A$ anwenden.
    \item \textbf{Kern}: L\"ose $Av = 0$. Freie Variablen $\to$ Basisvektoren des Kerns.
    \item \textbf{Bild}: Pivot-Spalten der Stufenform $\to$ entsprechende \emph{Original-Spalten} sind die Basis. \\
    \emph{Achtung}: Nicht die Spalten der Stufenform nehmen, sondern die der Original-Matrix!
    \item \textbf{Pr\"ufe Dimensionssatz}: $\dim \text{Ker} + \dim \text{Im} \stackrel{?}{=} \dim V$.
    \item \textbf{Inj./Surj./Bij.}: Aus den Dimensionen direkt ablesen.
\end{enumerate}
\end{merkbox}

\newpage

%=============================================================================
\part*{Teil 2: Aufgabe 35}

$T \colon \mathbb{R}^5 \to \mathbb{R}^3$. Schon wegen der Dimensionen kann $T$ niemals injektiv (also auch nie bijektiv/invertierbar) sein. Es bleibt nur die Frage, wie ``schlimm'' der Kern ist und ob $T$ wenigstens surjektiv ist.

\medskip

\textbf{Vorgehensweise:} Eine einzige Gau\ss-Elimination liefert die Antworten f\"ur (a), (b), (c). Pivots stehen bei den Variablen $a, b$ -- also sind $c, d, e$ freie Variablen im Kern, und die Spalten $1, 2$ von $A$ bilden die Basis des Bildes.

\medskip

\textbf{Konkret:} $\dim \text{Ker} = 3$, $\dim \text{Im} = 2$. Dimensionssatz: $3 + 2 = 5 = \dim \mathbb{R}^5$ \checkmark. \emph{Nicht injektiv} (Kern $\neq \{0\}$), \emph{nicht surjektiv} (Bild $\neq \mathbb{R}^3$), \emph{nicht invertierbar}.

\medskip

\textbf{(c) -- der Trick:} $T(v) = (1, 2, 5)$ l\"osen $=$ partikul\"are L\"osung $v_0$ finden, dann gesamten Kern dazu addieren. (Wie ``allgemeine L\"osung $=$ spezielle L\"osung $+$ homogene L\"osung'' beim Differentialgleichungen.)

\newpage

%=============================================================================
\part*{Teil 3: Aufgabe 36}

$T \colon \mathbb{R}^3 \to \mathbb{R}^3$. Jetzt Quadrat -- $T$ k\"onnte invertierbar sein.

\subsection*{Strategie f\"ur (a)}

Statt Matrix aufstellen und Determinante rechnen: direkt versuchen, $T^{-1}$ zu finden. Setze $(x, y, z) = T(a, b, c)$ und l\"ose nach $a, b, c$ auf. Wenn das eindeutig geht $\Rightarrow$ $T$ invertierbar; die L\"osung \emph{ist} dann $T^{-1}$.

\textbf{Warum funktioniert das?} ``Eindeutige R\"uckaufl\"osung'' bedeutet genau: zu jedem Output gibt es genau einen Input. Das ist bijektiv.

\subsection*{(b) und (c) folgen direkt aus (a)}

\begin{merkbox}[title={Faustregel}]
Wenn man wei\ss, dass $T$ invertierbar (also bijektiv) ist, dann braucht man \textbf{nichts mehr rechnen}:
\[
\text{Ker}(T) = \{0\}, \qquad \text{Im}(T) = W = \text{ganzer Zielraum}.
\]
Und $T$ ist automatisch injektiv, surjektiv, bijektiv.
\end{merkbox}

\newpage

%=============================================================================
\part*{Teil 4: Aufgabe 37}

$T \colon \mathbb{R}^{2\times 2} \to \mathbb{R}^{2\times 2}$, definiert als ``links mit $M$ multiplizieren''. Beachte: $\dim \mathbb{R}^{2\times 2} = 4$.

\subsection*{Schl\"usselbeobachtung}

Die Matrix $M = \begin{pmatrix} 1 & 2 \\ 4 & 8 \end{pmatrix}$ ist \textbf{singul\"ar} (Rang $1$: Zeile 2 $= 4 \cdot$ Zeile 1, Spalte 2 $= 2 \cdot$ Spalte 1).

\begin{merkbox}[title={Konsequenz}]
F\"ur $T(X) = M X$ mit Rang-$1$-Matrix $M$:
\begin{itemize}[nosep]
    \item Das Ergebnis $M X$ hat in jeder Spalte ein Vielfaches von Spalte 1 von $M$ (also $(1, 4)^\top$).
    \item Daher ist $\dim \text{Im}(T) = 2$ (je eine Dimension pro Spalte des Outputs).
\end{itemize}
\end{merkbox}

\subsection*{Direkter Weg}

$T(X)$ mit allgemeinem $X = \begin{pmatrix} a & b \\ c & d \end{pmatrix}$ ausrechnen, dann Bedingungen f\"ur Kern und Bild ablesen. Heraus kommt:
\[
T(X) = \begin{pmatrix} a + 2c & b + 2d \\ 4(a + 2c) & 4(b + 2d) \end{pmatrix}.
\]

Daraus direkt:
\begin{itemize}[nosep]
    \item Kern: $a + 2c = 0$ und $b + 2d = 0$. Zwei freie Variablen ($c, d$). $\dim \text{Ker} = 2$.
    \item Bild: jede Matrix der Form ``erste Zeile beliebig, zweite Zeile $= 4 \cdot$ erste''. $\dim \text{Im} = 2$.
\end{itemize}

Dimensionssatz: $2 + 2 = 4$ \checkmark.

\newpage

%=============================================================================
\part*{Teil 5: Aufgabe 38 -- Existenzfragen}

\begin{merkbox}[title={Hauptwerkzeug f\"ur Existenzaufgaben}]
\textbf{Schritt 1: Dimensionscheck.} Berechne $\dim \text{Ker}$ und $\dim \text{Im}$ \"uber den Dimensionssatz. Pr\"ufe, ob die Werte
\begin{enumerate}[nosep]
    \item \emph{ganze Zahlen} sind,
    \item $\geq 0$ sind,
    \item $\dim \text{Im} \leq \dim W$ (Zielraum) erf\"ullen.
\end{enumerate}
\textbf{Wenn ein Check fehlschl\"agt $\to$ Es existiert kein solches $T$.} \\
\textbf{Wenn alle Checks bestehen $\to$ Konstruiere ein konkretes Beispiel.}
\end{merkbox}

\subsection*{(a) $\mathbb{R}^2 \to \mathbb{R}^2$, $\text{Ker} = \text{Im}$}

$\text{Ker} = \text{Im}$ erzwingt gleiche Dimensionen. Dimensionssatz: $2 \dim \text{Ker} = 2 \Rightarrow \dim \text{Ker} = 1$. Ganzzahl, alles okay $\Rightarrow$ existiert.

Beispiel: $T(a, b) = (b, 0)$ ``vergisst $a$, schiebt $b$ in die erste Komponente''. Kern $=$ Bild $=$ $x$-Achse.

\subsection*{(b) $\mathbb{R}^3 \to \mathbb{R}^3$, $\text{Ker} = \text{Im}$}

$2 \dim \text{Ker} = 3 \Rightarrow \dim \text{Ker} = \tfrac{3}{2}$. \textbf{Keine ganze Zahl} $\Rightarrow$ existiert nicht.

\subsection*{(c) $\mathbb{R}^6 \to \mathbb{R}^3$, $\dim \text{Ker} = 2$}

Dimensionssatz: $\dim \text{Im} = 6 - 2 = 4$. Aber $\text{Im} \subseteq \mathbb{R}^3$, also $\dim \text{Im} \leq 3$. \textbf{Widerspruch} $\Rightarrow$ existiert nicht.

\begin{merkbox}[title={Muster erkennen}]
Solche Aufgaben sind meistens reine Dimensionsrechnungen -- die linearen Abbildungen sind nur ``Verpackung''. Schreibe die drei Zahlen $\dim V$, $\dim W$, $\dim \text{Ker}$ hin, rechne $\dim \text{Im} = \dim V - \dim \text{Ker}$ aus, und pr\"ufe, ob das in den Zielraum passt.
\end{merkbox}

\end{document}
